\scanpage{354}
% @ head of sheet
\mspage{78}
% @ top, lines 1–3
couples $(X_i\subset X,Y_i\subset Y)$ de sous-catégories de $X$ et de $Y$, telles que l'on ait $X_i,Y_i\in\mathrm{Ob}\,\mathrm{Cat}_{\unsure{\kappa|\pi}}$ (i.e. $\operatorname{card}(X_i),\operatorname{card}(Y_i)\leq\kappa$), et que $f(\deleted{\unsure{X}}X_i)\subset Y_i$.
% @ upper middle, lines 4–6
Il est immédiat que $\deleted{\unsure{\pi}}I$ est \unsure{filtrante|filtrant}, pour la relation d'ordre évidente
\[
(X_i,Y_i)\subset(X_j,Y_j)\quad\Longleftarrow\quad X_i\subset Y_i\text{ et }X_j\subset Y_j,
\]
% @ centre, lines 7–12
et $=$ \unsure{grde|grande} devant $\kappa$, que $I$ est stable par $\underrightarrow{\lim}$ filtrantes de cardinal $\leq\kappa$, et que le foncteur
\[
i\longmapsto(f_i:X_i\unsure{\mathbin{-}|\to}Y_i):I\longrightarrow\underline{\mathrm{Fl}}\,\mathrm{Cat}
\]
commute aux dites $\underrightarrow{\lim}$, enfin les $f_i$ sont dans $\underline{\mathrm{Fl}}(\mathrm{Cat})_\kappa$ par construction.
% @ lower middle, lines 13–15
\deleted{Il est \unsure{évidt|évident|clair}} Ainsi on a un système inductif $L_1$-$\kappa$-adapté, et \inserted{il est clair que} sa limite inductive est can. is. à $f$. \unsure{Ceci vu}

% @ lower third, corollary and cancellations
\inserted{\uuline{Corollaire} 5.3.2} \deleted{\unsure{vu}, le fait que} La condition : que $W$ soit $\pi$-accessible signifie donc que pour \unsure{toute} $f\inserted{{:}\,X\to Y}\in W$, \deleted{\unsure{si on l'ens} \unsure{ordonné}} et pour \unsure{tte} sous-catégorie $\unsure{X_i|X'}$ de $X$ et $Y_i$ de $Y$, de cardinal $\leq\pi$ (avec $f(X_i)\subset Y_i$, mais en fait c'est inutile, quitte à agrandir $\unsure{Y_i|X_i,Y_i}$), il existe \unsure{alors|deux} sous-catégories
% archive sheet number: 353
